\documentclass[lettersize,journal]{IEEEtran}
\usepackage{amsmath,amsfonts}
\usepackage{algorithmic}
\usepackage{algorithm}
\usepackage{array}
\usepackage[caption=false,font=normalsize,labelfont=sf,textfont=sf]{subfig}
\usepackage{textcomp}
\usepackage{stfloats}
\usepackage{url}
\usepackage{verbatim}
\usepackage{graphicx}
\usepackage{cite}
\usepackage{lipsum}
\hyphenation{op-tical net-works semi-conduc-tor IEEE-Xplore}
% updated with editorial comments 8/9/2021

\begin{document}

\title{Res-HMVCL: A Residual Heterogeneous Multi-View Contrastive Learning for Encrypted Traffic Classification under Extreme Label Scarcity}

\author{Hanif Nur Ilham Sanjaya , Alfan Presekal}

% The paper headers
\markboth{IEE Transaction Journal \LaTeX\ Class Files,~Vol.~xx, No.~xx, January~2026}%
{Shell \MakeLowercase{\textit{et al.}}: A Sample Article Using IEEEtran.cls for IEEE Journals}

\IEEEpubid{0000--0000/00\$00.00~\copyright~2021 IEEE}
% Remember, if you use this you must call \IEEEpubidadjcol in the second
% column for its text to clear the IEEEpubid mark.

\maketitle

\begin{abstract}
 The rapid proliferation of encryption protocols (SSL/TLS) has rendered Deep Packet Inspection (DPI) ineffective, necessitating flow-based Deep Learning (DL) classifiers. However, existing DL approaches face degradation under label scarcity and struggle to distinguish complex streaming behaviors. To address these limitations, we developed \textbf{Res-HMVCL}, a framework integrating 1D-CNN payload representations, bi-directional flow statistics, and Spectral (FFT) features via a residual contrastive mechanism. We introduce \textit{Universal Contrastive Pre-training} to learn traffic manifolds from unlabeled data. Experimental results on ISCXVPN2016 demonstrate an F1-Score of 94\% and 87\% for binary and category classification respectively using only 20\% labeled data, with a critical stability threshold identified at 10\% labeled data, significantly outperforming SOTA baselines in data efficiency.
\end{abstract}

\begin{IEEEkeywords}
Cyber security, Digital Substation, Kyber, DIlithium , GOOSE, HMAC, Post-Quantum Cryptography .
\end{IEEEkeywords}

\section{Introduction}
\IEEEPARstart{T}{he} rapid evolution of internet privacy mechanisms has led to the ubiquitous adoption of encryption protocols. As noted by Kahn \cite{kahn1996codebreakers}, encryption techniques have advanced significantly to meet the demand for secure communication across military, commercial, and personal domains. Among these technologies, Virtual Private Networks (VPNs) have become a fundamental tool for users seeking to bypass geo-restrictions, secure remote access, and safeguard sensitive data from unauthorized surveillance \cite{singhpall1996vpn}. However, while VPNs provide essential privacy benefits, they also introduce significant challenges for network management and security monitoring. The anonymity provided by VPNs can be exploited for malicious activities, creating a complex "dual-use" dilemma where legitimate privacy preservation conflicts with the need to detect cyber threats \cite{dipietro2021new}.

Traditionally, network traffic classification relied on port-based methods, which matched packet headers to the Internet Assigned Numbers Authority (IANA) registry \cite{yoon2009internet}. However, empirical studies have shown that port-based classification is increasingly unreliable, with failure rates ranging from 30\% to 70\% due to the use of dynamic port negotiation and obfuscation techniques \cite{moore2005toward}. To address this, Deep Packet Inspection (DPI) emerged as a robust alternative, analyzing packet payloads for specific signatures \cite{finsterbusch2013survey}. Yet, the widespread deployment of modern encryption standards renders DPI largely ineffective, as payload content is obfuscated and unreadable without decryption keys \cite{wang2017endtoend}.

Consequently, the research community has shifted toward Machine Learning (ML) and Deep Learning (DL) paradigms. Recent surveys indicate a clear trend moving from manual feature engineering toward DL models that utilize raw packet data or time-series representations \cite{razooqi2025vpn}. Approaches such as "Deep Packet" \cite{lotfollahi2019deep} and 1D-CNNs \cite{wang2017endtoend} have demonstrated state-of-the-art performance by automatically learning complex patterns from encrypted flows. However, these supervised DL models face two critical limitations. First, they require massive amounts of labeled data, which is labor-intensive and expensive to acquire in real-world environments \cite{menezes2018handbook}. Second, existing semi-supervised approaches \cite{lin2022semisupervised} often fail to effectively align the multi-modal nature of network traffic. They typically treat payload signatures and statistical flow behaviors as separate entities, missing the opportunity to use one view to semantically guide the representation learning of the other.

To bridge this gap, we propose LE-MVCL (Label-Efficient Multi-View Contrastive Learning), a novel framework designed to achieve robust encrypted traffic classification under severe label scarcity. LE-MVCL distinguishes itself through an \textit{asymmetric} pre-training strategy that leverages unlabeled data to align a Deep 1D-CNN (Payload View) with a Deep MLP (Statistical View). Furthermore, we introduce a comprehensive feature synthesis encompassing payload, flow dynamics, and burst behaviors, ensuring that granular statistical precision is preserved alongside deep semantic embeddings.

The main contributions of this paper are as follows:
\begin{itemize}
    \item The LE-MVCL Framework: We propose a novel framework that adapts Contrastive Learning to the domain of VPN identification. Unlike standard supervised methods, our approach leverages unlabeled data to learn robust flow representations, achieving 94.5\% Binary F1 with only 10\% labeled data and 96.0\% with 30\% labeled data, significantly mitigating the label scarcity problem.
    
    \item Asymmetric Multi-View Pre-training: We introduce an asymmetric contrastive pre-training strategy that aligns a Deep 1D-CNN (Payload View) with a Deep MLP (Statistical View). We demonstrate that using the statistical view as a "semantic guide" during pre-training forces the CNN to learn high-level behavioral patterns from raw bytes, significantly outperforming standard end-to-end training.
    
    \item Hybrid Late-Fusion Architecture: We empirically demonstrate that a Hybrid Late-Fusion strategy—combining learned semantic embeddings with raw statistical features—yields superior performance compared to pure deep learning approaches. This validates that preserving the granular precision of statistical metrics is critical for distinguishing fine-grained VPN applications (e.g., Skype vs. Hangouts).
    
    \item Robust Feature Engineering: We synthesize a comprehensive multi-view feature set comprising the Alpha (Payload), Beta (Flow Dynamics), Gamma (Burst Behavior), and FFT (Spectral) components. We prove that fusing these diverse views enables the detection of obfuscated traffic patterns, maintaining high efficacy ($>$93\% App F1) even when distinguishing between isotocol services.
\end{itemize}
\input{section2_design}


\section{Conclusion}
The conclusion goes here \cite{ref1}.


\section*{Acknowledgments}
This should be a simple paragraph before the References to thank those individuals and institutions who have supported your work on this article.







%{\appendices
%\section*{Proof of the First Zonklar Equation}
%Appendix one text goes here.
% You can choose not to have a title for an appendix if you want by leaving the argument blank
%\section*{Proof of the Second Zonklar Equation}
%Appendix two text goes here.}



% References
\bibliographystyle{IEEEtran}
\bibliography{references}


\newpage

\section{Biography Section}
If you have an EPS/PDF photo (graphicx package needed), extra braces are
 needed around the contents of the optional argument to biography to prevent
 the LaTeX parser from getting confused when it sees the complicated
 $\backslash${\tt{includegraphics}} command within an optional argument. (You can create
 your own custom macro containing the $\backslash${\tt{includegraphics}} command to make things
 simpler here.)
 
\vspace{11pt}

\bf{If you include a photo:}\vspace{-33pt}
\begin{IEEEbiography}[{\includegraphics[width=1in,height=1.25in,clip,keepaspectratio]{fig1}}]{Michael Shell}
Use $\backslash${\tt{begin\{IEEEbiography\}}} and then for the 1st argument use $\backslash${\tt{includegraphics}} to declare and link the author photo.
Use the author name as the 3rd argument followed by the biography text.
\end{IEEEbiography}

\vspace{11pt}

\bf{If you will not include a photo:}\vspace{-33pt}
\begin{IEEEbiographynophoto}{John Doe}
Use $\backslash${\tt{begin\{IEEEbiographynophoto\}}} and the author name as the argument followed by the biography text.
\end{IEEEbiographynophoto}




\vfill

\end{document}

